\documentclass[11pt,a4paper]{article}

\usepackage[utf8]{inputenc}

\usepackage[T1]{fontenc}

\usepackage[polish]{babel}

\usepackage{lmodern}

\usepackage{geometry}

\geometry{a4paper,margin=27mm}

\usepackage{microtype}

\usepackage[hidelinks]{hyperref}

% Portfolio: MSC-AR-P26; source: prospective ARC-XR research scope.

\title{Faza ruchu dłoni a optymalny moment wyświetlania instrukcji w AR}

\author{Mateusz Pomianek}

\date{Wersja 0.3.0, 21 listopada 2026}

\begin{document}

\maketitle

\begin{abstract}

Przedmiotem opracowania jest Faza ruchu dłoni a optymalny moment wyświetlania instrukcji w AR. Czy moment wyświetlenia wskazówki przed ruchem powoduje mniej błędów niż podczas ruchu? W neutralnej scenie AR porównane zostaną wskazówki wyświetlane przed rozpoczęciem ruchu i podczas jego wykonywania, przy identycznej treści instrukcji. Głównym punktem końcowym jest odsetek błędnych kroków sekwencji manualnej. Opis ma charakter wstępny i nie przedstawia niezweryfikowanych wyników.

\end{abstract}

\section{Wprowadzenie}

Rozwój adaptacyjnych systemów XR stwarza możliwość badania interakcji człowieka z informacją przestrzenną, ale samo zastosowanie AR, VR lub modeli sztucznej inteligencji nie przesądza o przewadze poznawczej czy ergonomicznej. W rozpatrywanym zagadnieniu, Faza ruchu dłoni a optymalny moment wyświetlania instrukcji w AR, kluczowe jest rozdzielenie właściwości systemu od zmienności zachowania użytkownika. Porównanie należy zaprojektować tak, aby mierzyć jednoznaczne zjawisko, a nie tylko ogólną ocenę technologii.

\section{Cel pracy i pytanie badawcze}

Czy moment wyświetlenia wskazówki przed ruchem powoduje mniej błędów niż podczas ruchu? Projekt wymaga krytycznego przeglądu aktualnego stanu badań oraz zdefiniowania nowości w stosunku do innych tematów portfolio XR.

\section{Zakres i odrębny wkład naukowy}
Znaczenie fazy wykonania ruchu dla czasu podania instrukcji AR przy identycznej treści, odróżnione od szybkości prezentacji instrukcji i korekty gestu. Nowość musi dotyczyć efektu fazy ruchu na poprawność procedury, nie twierdzenia, że sama synchronizacja AR jest niespotykana.
\section{Hipotezy badawcze}
\paragraph{H1.} Podanie tej samej wskazówki przed rozpoczęciem ruchu powoduje mniej błędów kroków niż prezentacja w jego trakcie.
\paragraph{H2.} Efekt jest największy w krokach wymagających precyzji chwytu i integracji wzrokowo-ruchowej, nie w prostych krokach.
\paragraph{H3.} Zysk pozostaje po uwzględnieniu faktycznego czasu dostarczenia instrukcji i błędu detekcji początku ruchu.
\section{Projekt badania i plan eksperymentów}
\subsection{A. Operacjonalizacja i kwalifikacja danych}
Znaczenie fazy wykonania ruchu dla czasu podania instrukcji AR przy identycznej treści, odróżnione od szybkości prezentacji instrukcji i korekty gestu. Jednostka wnioskowania: uczestnik × krok sekwencji manualnej z niezależną fazą ruchu. Przed badaniem należy określić źródła, wersje urządzeń, czas akwizycji, etykiety i niezależną referencję. Nie liczyć zależnych okien i klatek jako niezależnych prób.
\subsection{B. Eksperyment porównawczy}
Przygotować sekwencje chwytu z niezależnie rejestrowaną fazą ruchu; warunki wskazówka przed rozpoczęciem i w trakcie, z tym samym czasem widoczności. Wynik pierwszorzędowy: różnica odsetka błędnych kroków: wskazówka przed ruchem vs wskazówka w trakcie ruchu. Zamrozić losowanie warunków, progi, mianownik punktu końcowego i kryteria wyłączeń bez dostępu do oceny końcowej.
\subsection{C. Kontrole i analiza odporności}
Kalibrować wideo/kinematykę i czas renderowania, ustalić reguły nieklasyfikowalnych ruchów i nie wybierać wygodniejszego momentu po obserwacji błędów. Zbadać przypadki błędnego wykrycia momentu rozpoczęcia ruchu, bez wnioskowania o klinicznej rehabilitacji. Utrata sygnału, abstencja, przerwanie zadania i dane brakujące są odrębnymi kategoriami wyniku.
\subsection{D. Walidacja na niezależnym materiale}
Walidacja musi zachować podział po uczestniku, sesji, scenie lub zadaniu. Model logistyczny błędu kroku z klastrowaniem po uczestniku i typie chwytu; analiza latencji jako czynnik zakłócający. Raportować wyniki negatywne i rozbieżności z założeniami.
\section{Estymand, analiza statystyczna i decyzje}
różnica odsetka błędnych kroków: wskazówka przed ruchem vs wskazówka w trakcie ruchu. Jednostka: uczestnik × krok sekwencji manualnej z niezależną fazą ruchu. Model logistyczny błędu kroku z klastrowaniem po uczestniku i typie chwytu; analiza latencji jako czynnik zakłócający. H1 jest planowanym kontrastem głównym; H2 i H3 pozostają pomocnicze do zatwierdzenia hierarchii testów oraz kontroli wielokrotności. Wartości marginesów, mocy, minimalnego efektu istotnego praktycznie i liczebności mają wynikać z pilotażu, a nie z post hoc wyboru wyniku.
\section{Stan badań i granica nowości}
\begin{itemize}
\item AR guidance timing has already been manipulated experimentally; novelty needs phase-specific step errors, not generic scheduling. \cite{xrprior1}.
\item Prior hand-gesture AR guidance has tested error feedback; establish a distinct motor-phase or procedural endpoint. \cite{xrprior2}.
\end{itemize}
Nowość musi dotyczyć efektu fazy ruchu na poprawność procedury, nie twierdzenia, że sama synchronizacja AR jest niespotykana. Odrębność na tle opublikowanych metod i pozostałych artykułów wymaga osobnego audytu.
\section{Warunki bezpieczeństwa, sprzętu i zgodności}
To prospektywna koncepcja, nie zgoda na eksperyment, prerejestracja lub wynik empiryczny. Praca z uczestnikami, interwencje stresowe i pomiar biometryczny wymagają przeglądu etycznego, świadomej zgody i ochrony danych. Quest 3 w wersji standardowej nie posiada eye trackingu, a Tobii Eye Tracker 5 jest urządzeniem ekranowym; nie zastępować jednego stanowiska drugim.
\section{Materiał i metody}

W neutralnej scenie AR porównane zostaną wskazówki wyświetlane przed rozpoczęciem ruchu i podczas jego wykonywania, przy identycznej treści instrukcji. Wymagane jest zdefiniowanie jednostki analizy, zasad doboru przypadków, wersji scenariusza i źródeł pomiaru. Rozkład warunków oraz reguły wykluczenia prób należy określić przed pomiarem. W badaniach z udziałem ludzi niezbędna jest uprzednia właściwa ocena etyczna i świadoma zgoda. W testach wyłącznie symulacyjnych należy oddzielić prawdę symulacyjną od obserwacji empirycznej.

\subsection{Wynik główny i kryteria jakości}

Odsetek błędnych kroków sekwencji manualnej. Pomiar powinien zachowywać jawne jednostki, znaczniki czasu i reguły oceny poprawności. Przy modelach predykcyjnych stosować rozłączne partycje oceny; przy eksperymentach interwencyjnych oddzielić skuteczność predykcji od przyczynowego wpływu interfejsu.

\section{Plan analizy}

Model mieszany dla prób i uczestników, osobno ocenić czas reakcji i faktyczną latencję renderowania. W wynikach należy raportować wielkość efektu lub błąd estymacji z przedziałem niepewności, liczbą obserwacji i odsetkiem odrzuconych prób. Analizy dodatkowe wymagają jasnego oznaczenia jako eksploracyjne. Nie wybierać progów decyzyjnych ani modelu na podstawie obserwacji ze zbioru testowego.

\section{Wyniki}

W obecnej wersji nie raportuje się wyników pomiarów, tabel statystycznych ani parametrów wyznaczonych z nieudokumentowanych danych. To kierunek badawczy, a nie zatwierdzony protokół badania ani zgłoszony artykuł.

\section{Dyskusja}

Interpretacja przyszłych wyników musi uwzględnić alternatywne mechanizmy i zależność wniosków od jakości danych, warunków zadania oraz ograniczeń sprzętu. Wartość potencjalnego rozwiązania zależy nie tylko od trafności modelu, lecz także od jego opóźnienia, stabilności i wpływu na zachowanie człowieka. Wyniki testów syntetycznych nie mogą być automatycznie uogólniane na populację uczestników.

\section{Ograniczenia badania}

Błędne wykrycie rozpoczęcia ruchu może zafałszować kontrast; temat pozostaje koncepcją do audytu nowości. Granice nowości i trafności ekologicznej wymagają niezależnej weryfikacji. Żadna część szkicu nie dowodzi przeprowadzenia zbierania danych, uzyskania zgód etycznych ani gotowości manuskryptu do zgłoszenia.

\section{Wnioski}

Przedstawiono problem, kontrast i sposób jego planowanej oceny. Rozstrzygnięcie postawionego pytania wymaga przygotowania protokołu, zweryfikowanych źródeł oraz poprawnie udokumentowanych wyników. Bez tych materiałów nie można sformułować konkluzji empirycznych.

\begin{thebibliography}{9}
\bibitem{xrprior1} The Art of Timing: Effects of AR Guidance Timing on Speed Control (2024), ISMAR 2024. \url{https://www.krisluyten.net/publications/ceyssens2024/}.
\bibitem{xrprior2} Exploring AR hand augmentations as error feedback mechanisms for enhancing gesture-based tutorials (2025), Frontiers in Virtual Reality. \url{https://doi.org/10.3389/frvir.2025.1574965}.
\end{thebibliography}
\end{document}

